\begingroup\color{blue}
% =====================================================================
%  Branch and Price Solution Approach for 3D-SCPTOP
% =====================================================================
\section{Branch and Price Solution Approach}
\label{sec:math}

There are three structural difficulties with the 3D-SCPTOP model of
Section \ref{sec:model}. The first is that a priori we do not know how
many loading configurations are needed, so $|L|$ must be
overestimated; every configuration that turns out to be unnecessary
still carries its full block of placement variables. The second is
that the model replicates a complete two-dimensional packing problem
$|G|\cdot|L|$ times, and each replica contributes
$4\binom{|I|}{2}$ big-$M$ disjunctions
\eqref{sb:eq:nonov-x1}--\eqref{sb:eq:nonov-y2} whose linear
relaxation is weak. The third is that the configurations are mutually
interchangeable, so the search tree is massively symmetric.

It would overstate the case to call the resulting model intractable in
general, and Section \ref{sec:results} shows why: with a good
incumbent supplied, instances spanning five periods are solved to
proven optimality in seconds, even at sixty thousand variables. What
defeats it is the planning horizon. The placement block is replicated
per configuration rather than per period, but the truck-assignment and
inventory blocks grow with $|T|$, and it is their interaction with the
symmetric configuration design that explodes: at ten periods the model
returns no feasible solution within half an hour. Operational
procurement planning lives at horizons of weeks, so that is the regime
the method has to serve, and it is the regime in which a decomposition
earns its place. A column generation scheme based on branch and price
is therefore developed.

The key components of the branch and price algorithm are the master
problem (MP), the restricted master problem (RMP), the pricing
problem (PP) and the branching policy. The MP is a reformulation of
the model that exploits its decomposable structure. Inspecting
Section \ref{sec:model}, every geometric constraint
\eqref{sb:eq:slot-active}--\eqref{sb:eq:util} carries a single
configuration index $l$, and every loading constraint
\eqref{sb:eq:slot-cap}--\eqref{sb:eq:stack-w-sup} carries a single
truck index $k$ and a single period $t$; the transport and inventory
block sees them \emph{only} through the aggregate flows $Q_{nt}$ and
$Q_{rt}$ of \eqref{sb:eq:flow-n}--\eqref{sb:eq:flow-r}. The model is
therefore block-angular, with one independent block per dispatched
truck. This suggests taking as the unit of decision neither the
individual RTI nor the loading configuration, but the
\emph{complete loaded truck}.

The MP is typically too large to be solved directly because it
involves a huge number of columns representing all possible truck
loads. To manage this complexity we consider only a smaller subset of
columns, which leads to the RMP. Solving the RMP provides an optimum
for the included subset, but not necessarily for the entire column
set. RMP integrality constraints are therefore relaxed to obtain dual
variables, and the PP uses those duals to identify columns of
negative reduced cost that should be added to the RMP. If no such
column exists, the integrality-relaxed MP has been solved to
optimality. If all the variables associated with the solution are
integers, an integer solution to the current branch has been found.
If not, the branch must be divided by adding further constraints.
Figure \ref{fig:B&P3D} exemplifies the algorithm.

\begin{figure}[h!]
    \centering
    \resizebox{0.8\columnwidth}{!}{%
\begin{tikzpicture}[node distance=2cm, auto, thick]
    \tikzstyle{block} = [rectangle, draw=white, fill=white, minimum height=2em,
                         text width=17em, text centered]
    \tikzstyle{line} = [draw, -Latex]

    \node [block] (MP) {\textbf{Master Problem}\\[2pt]
        \footnotesize all feasible truck loads $P^{p}\cup P^{s}$};
    \node[block, below=0.9cm of MP] (RMP) {\textbf{Restricted Master Problem}\\[2pt]
        \footnotesize loads generated so far, $\lambda\ge 0$};
    \node[block, below=0.9cm of RMP] (PP) {\textbf{Pricing Problem}\\[2pt]
        \footnotesize design one truck: 2D placement,\\
        \footnotesize layering and product assignment};

  \draw[bend right,left,->,align=left] (MP.west) to
        node {A reduced number \\ of columns} (RMP.west);
  \draw[bend right,left,->,align=left] (RMP.west) to
        node {Dual prices $\pi_{ct}$, $\nu_{gt}$} (PP.west);
  \draw[bend left,right,<-,align=left] (RMP.east) to
        node {If reduced cost is negative \\ add the truck load as a column; \\
              if not, branch} (PP.east);
  \draw[bend right,below,->,align=center] (PP.west) to
        node {Maximise $\sum_{c}\pi_{ct}a_{c}$ \\ subject to
              \eqref{eq:PPtype}--\eqref{eq:PPsymm}} (PP.east);
  \draw[bend left,below,<-,align=center] (RMP.west) to node {} (RMP.east);
\end{tikzpicture}
}
\caption{Branch and price scheme for 3D-SCPTOP. Every geometric,
height, crush-limit, payload and full-truckload constraint is enforced
inside the pricing problem, so each column is a physically loadable
truck.}
\label{fig:B&P3D}
\end{figure}


\subsection{Formulation of the Master Problem and the Restricted
Master Problem}
\label{sec:mp}

The key idea of the set-partitioning formulation for 3D-SCPTOP is to
enumerate all the feasible \emph{truck loads} with their associated
costs. A plant-bound truck load is a vector
$\mathbf{a}=(a_{n})_{n\in N}$ of full RTIs for which there exist a
slot count $m\le|I|$, stack types $h_{1},\dots,h_{m}\in H$,
orientations, bottom-left coordinates $(x_{i},y_{i})$, layer counts
$\lambda_{ir}$ and per-slot product loads $v_{in}$ such that
\begin{align}
&\text{the $m$ footprints, rotated as chosen, do not overlap and lie
   inside } pK\times qK,
   \label{eq:col-geom}\\
&\textstyle\sum_{r\in R(h_{i})} r_{pr}\lambda_{ir}\;\le\;rK
   &&\forall i\le m,
   \label{eq:col-height}\\
&\textstyle\sum_{n\in M(r)} v_{in}\;=\;n_{h_{i}r}\lambda_{ir}
   &&\forall i\le m,\;r\in R(h_{i}),
   \label{eq:col-assign}\\
&\textstyle\sum_{n\in N} w_{n}v_{in}\;\le\;W^{\max}
   &&\forall i\le m,
   \label{eq:col-crush}\\
&\textstyle\sum_{n\in N} w_{n}a_{n}\;\le\;pK^{w},
   \qquad
   \max\Bigl\{\tfrac{\sum_{n}\gamma_{p\rho(n)}a_{n}}{\Gamma_{K}},\;
              \tfrac{\sum_{n}w_{n}a_{n}}{pK^{w}}\Bigr\}\;\ge\;\alpha,
   \label{eq:col-truck}\\
&a_{n}=\textstyle\sum_{i\le m} v_{in}
   &&\forall n\in N.
   \label{eq:col-agg}
\end{align}
Let $P^{p}$ be the set of all such loads, and $P^{s}$ the analogous
set of supplier-bound loads $\mathbf{b}=(b_{r})_{r\in R}$, built with
folded heights $r_{sr}$, folded weights $w^{\mathrm{f}}_{r}$ and
folded volumes $\gamma_{sr}$; there no product assignment is needed,
because the weight of a folded RTI depends only on its type.

Two features of this definition matter. First, a column carries the
\emph{per-slot} product assignment $v_{in}$, so the crush limit
\eqref{eq:col-crush} is evaluated on the product mix actually riding
in that slot. The column therefore reproduces the period-exact cap
\eqref{sb:eq:stack-w-plant} of the compact model, and not the
conservative surrogate $\bar w_{r}=\max_{n\in M(r)}w_{n}$ that a pool
of capacity-only configurations is forced to use. Second, a column is
a feasibility certificate: the placement
$(h_{i},x_{i},y_{i},\text{rotation})$ witnessing
$\mathbf{a}\in P^{p}$ is retained, so any solution can be handed to a
loading crew as a drawing and re-verified independently of the solver
that produced it.

Constraint \eqref{eq:col-truck} is the same disjunction as
\eqref{sb:eq:util}--\eqref{sb:eq:utilw}, written on the load a truck
actually carries rather than on the capacity its configuration
provides. Section \ref{sec:results} shows that both regimes occur
across the benchmark: at $v=60$ the trailer fills by volume, at
$v=120$ it fills by weight.

Let $\lambda^{g}_{pt}\in\mathbb{Z}_{\ge0}$ count the trucks
dispatched in direction $g$ and period $t$ carrying load $p$, and let
$a^{p}_{n}$ and $a^{p}_{r}$ be the RTIs of product $n$, respectively
of type $r$, that load $p$ contains. Every truck costs $cK$. The MP
of our problem is
\begin{align}
    \min\;\; cK \sum_{g\in G}\sum_{t\in T}\sum_{p\in P^{g}}\lambda^{g}_{pt}
    \;+\;\sum_{g\in G}\sum_{n\in N}\sum_{t\in T} ic_{gn} I_{gnt}
    \label{eq:MPObjective}
\end{align}
subject to
\begin{align}
    \sum_{p\in P^{p}} a^{p}_{n}\lambda^{p}_{pt} = Q_{nt}
      \quad & \forall n\in N,\;t\in T \label{eq:MPlinkp}\\
    \sum_{p\in P^{s}} a^{p}_{r}\lambda^{s}_{pt} = Q_{rt}
      \quad & \forall r\in R,\;t\in T \label{eq:MPlinks}\\
    \sum_{p\in P^{g}} \lambda^{g}_{pt} \le |K|
      \quad & \forall g\in G,\;t\in T \label{eq:MPfleet}\\
    \text{\eqref{sb:eq:inv-plant-main}--\eqref{sb:eq:safety}}
      \quad & \text{(inventory balances and coverage stock)}
      \nonumber\\
    \lambda^{g}_{pt}\in\mathbb{Z}_{\ge0}
      \quad & \forall g\in G,\;p\in P^{g},\;t\in T. \nonumber
\end{align}
Objective \eqref{eq:MPObjective} minimises transport plus inventory
holding cost. Constraints
\eqref{eq:MPlinkp}--\eqref{eq:MPlinks} state that the flow of full
and folded RTIs is exactly what the dispatched trucks carry, and
\eqref{eq:MPfleet} caps the fleet available per direction and period.

MP and 3D-SCPTOP have the same optimal value whenever $|L|$ is not
restricted a priori. Given a 3D-SCPTOP solution, each used truck
$(k,t)$ induces a load vector satisfying
\eqref{eq:col-geom}--\eqref{eq:col-agg}, obtained by reading the slot
data off its configuration and setting
$v_{in}=V^{\mathrm{slot}}_{kilnt}$; counting trucks by load vector
gives a MP solution of equal cost. Conversely, introducing one
configuration per distinct load used by a MP solution and assigning
trucks accordingly reconstructs a 3D-SCPTOP solution of the same
cost. This equivalence is what distinguishes the present approach
from a pool-based method such as PLC-SCPTOP: a predefined pool is a
\emph{restriction} and can only yield upper bounds of unknown
quality, whereas MP is the same problem written differently, so any
bound derived from it is a bound on the original.

The RMP replaces $P^{g}$ by a subset $\bar P^{g}$ and relaxes
integrality,
\begin{align}
    \min\;\; cK \sum_{g\in G}\sum_{t\in T}\sum_{p\in \bar P^{g}}\lambda^{g}_{pt}
    \;+\;\sum_{g\in G}\sum_{n\in N}\sum_{t\in T} ic_{gn} I_{gnt}
    \label{eq:RMPObjective}
\end{align}
subject to \eqref{eq:MPlinkp}--\eqref{eq:MPfleet} restricted to
$\bar P^{g}$, with $\lambda^{g}_{pt}\ge0$. Dual values are obtained
by solving the RMP. Let $\pi_{nt}$ and $\pi_{rt}$ be the duals of
\eqref{eq:MPlinkp}--\eqref{eq:MPlinks} and $\nu_{gt}\le0$ those of
\eqref{eq:MPfleet}. The reduced cost associated with a plant-bound
load $p\in P^{p}$ dispatched in period $t$ is
\begin{equation}
\bar c(p,t)\;=\;cK-\sum_{n\in N}\pi_{nt}a^{p}_{n}-\nu_{pt},
\label{eq:redcost}
\end{equation}
and analogously in the supplier direction. A load with negative
reduced cost can improve the current RMP solution. For each iteration
we solve the subproblem of finding the load of minimum reduced cost;
loads with negative reduced cost are added as new columns, the RMP is
re-solved to generate new duals, and the process repeats until no
further improving column is discovered.

\subsection{Lower bounds}
\label{sec:bounds}

Care is needed here, because the natural candidate is not a bound at
all. The RMP is a \emph{restriction} of the MP: it optimises over
$\bar P^{g}\subseteq P^{g}$, so
$z_{\mathrm{RMP}}\ge z_{\mathrm{MP}}$. Reporting the RMP value as a
lower bound is therefore valid only once column generation has
converged, and wrong before then. We use two bounds accordingly.

\paragraph{Certified bound.}
When the PP has \emph{proven} that
$\zeta_{gt}:=\max_{p\in P^{g}}\sum_{c}\pi_{ct}a^{p}_{c}\le cK-\nu_{gt}$
for every $g$ and $t$, no column outside $\bar P^{g}$ has negative
reduced cost, so the current duals are feasible for the dual of the
full MP. Weak duality then gives
$z_{\mathrm{RMP}}=z_{\mathrm{MP}}\le z^{\star}_{\text{3D-SCPTOP}}$.
We report such bounds as \emph{certified}; the gap against the final
integer solution is then a genuine optimality gap, and not a gap with
respect to a pool.

\paragraph{Lagrangian bound.}
Certification requires the two-dimensional packing subproblem to be
closed to proven optimality, which is not always possible inside a
sensible time limit. Dualising
\eqref{eq:MPlinkp}--\eqref{eq:MPlinks} with the current $\pi$ leaves a
subproblem in which each column enters with reduced cost
\eqref{eq:redcost}; since the number of trucks an optimal solution can
dispatch in total is bounded by $\bar X$, the correction
\begin{equation}
z_{\mathrm{MP}}\;\ge\;z_{\mathrm{RMP}}
   \;+\;\bar X\cdot\min_{g\in G,\,t\in T}\;
   \min\bigl\{0,\;cK-\zeta_{gt}-\nu_{gt}\bigr\}
\label{eq:lagbound}
\end{equation}
holds at every iteration, converged or not. Two choices make
\eqref{eq:lagbound} usable rather than vacuous. The value substituted
for $\zeta_{gt}$ is the solver's own dual bound on the PP, never its
incumbent, so a subproblem stopped early still yields a valid
verdict. And $\bar X$ must be a genuine bound on the fleet actually
used: the textbook form charges $|K|$ separately for every direction
and period, which here is roughly an order of magnitude too
pessimistic and drove the reported gaps above $80\%$. Because
transport cost alone cannot exceed the objective, any incumbent
$\bar z$ gives $\bar X=\bar z/cK$, which is both valid and tight
enough to be informative.

\paragraph{What we report.}
We report the certified bound, and when certification fails we report
no bound at all. The Lagrangian correction \eqref{eq:lagbound} is
computed and recorded, but not published. In our implementation it
repeatedly returned values above solutions we could exhibit: on one
instance it gave $21552.23$ while the pool of that same run admits a
slack-free plan, verified truck by truck, costing $21502.80$. What
makes this dangerous rather than merely loose is that no check inside
the run detects it --- the run's own restricted master lies above the
true optimum too, so the estimate looks consistent from the inside,
and only an externally exhibited plan exposes it. A quantity that
cannot be falsified from within the method is not one we are willing
to print as a bound.

\paragraph{Two ways a certificate can be void.}
A certificate is only as good as the exhaustiveness of the check
behind it, and two implementation details silently broke ours. First,
the sweep skipped any $(g,t)$ whose commodity duals were all
non-positive, on the argument that nothing there can pay for a
truck. The reduced cost is
$cK-\nu_{gt}-\sum_{c}\pi_{ct}a_{c}$: non-positive duals bound the sum
by zero, so the argument establishes only $rc\ge cK-\nu_{gt}$ and holds
just when $cK-\nu_{gt}\ge 0$. The test never looked at $\nu$. A pair
skipped this way contributes neither a proof nor an objection, so the
certificate was claimed over pairs nobody had examined --- and it is
the only path by which that can happen, since a subproblem stopped on
time already reports itself unproven. Second, when the pricing model
is infeasible under a threshold constraint, the oracle returned a
bound of zero. Infeasibility there proves only that no column beats
the threshold, so the \emph{threshold} bounds $\zeta_{gt}$; reporting
zero understates $\zeta_{gt}$ and inflates \eqref{eq:lagbound}. Both
defects push the bound the same way, upwards, which is where every
invalid bound we observed sat.

\paragraph{When certification is attempted.}
Certifying at the root is wasteful, because the tail of column
generation buys almost nothing. Over ten recorded runs, stopping once
the best reduced cost exceeded $-0.2$ gave up $0.00$ of LP value in
nine of them and $0.59$ in $18795$ in the tenth, while saving between
$62$ and $764$ seconds --- time that buys tree nodes, which is where
the columns that change the fleet actually come from. We therefore
stop on that tolerance and spend a single exact sweep at the end, over
the final pool. If nothing prices out there and every $(g,t)$ returned
a bound, the LP value of that pool is a valid lower bound. It is
weaker than an unconverged root value, since more columns can only
lower the LP; unlike it, it is true.

\paragraph{Rounding the duals.}
The prices $\pi$ are irrational in general and must be rounded before
entering an integer-coefficient subproblem. We round them
\emph{upwards}, so the rounded objective dominates the true one and
the returned bound remains an over-estimate of $\zeta_{gt}$, as both
bounds above require. Rounding to nearest --- the obvious choice ---
makes the bound a slight under-estimate and silently invalidates the
certificate; in our first implementation it returned $1000.1$ against
a threshold of $1000.0$ on instances where the reduced cost was
exactly zero, and convergence was never declared.

\subsection{Pricing Subproblem}
\label{sec:pp}

We solve the pricing problem with a constraint programming approach.
The indices and parameters are the same as in the model proposed
before, but the decision variables do not depend on the configuration
index. Let $c$ range over products when $g=p$ and over RTI types when
$g=s$, let $\rho(c)$ be the RTI type carrying $c$, and let $S$ be the
slot bound \eqref{sb:eq:Ibound}. Table
\ref{tab:NomenclaturePricing3D} describes the new decision variables.

\begin{longtable}[hbt!]{|p{0.10\textwidth}|p{0.84\textwidth}|}
\caption{Nomenclature of the pricing subproblem}
\label{tab:NomenclaturePricing3D}\\
\hline
\multicolumn{2}{|l|}{\textbf{Decision variables}}\\ \hline
$u_{i}$        & Binary variable indicating whether slot $i$ is used in the new load\\
$\chi_{ih}$    & Binary variable indicating whether slot $i$ holds a stack of type $h$\\
$\varrho_{ih}$ & Binary variable indicating whether that stack is rotated by $90^{\circ}$\\
$\delta^{x}_{i},\delta^{y}_{i}$ & Footprint of slot $i$ along the $x$ and $y$ axes\\
$x_{i},y_{i}$  & Bottom-left coordinates of slot $i$ on the truck floor\\
$\ell_{ir}$    & Number of layers of RTI type $r$ in slot $i$\\
$v_{ic}$       & Number of RTIs of commodity $c$ loaded in slot $i$\\
$a_{c}$        & Number of RTIs of commodity $c$ in the new load\\
$\theta$       & Binary variable indicating whether the full-truckload requirement is met on volume rather than on weight\\
\hline
\end{longtable}

The objective is to find new truck loads with negative reduced cost.
Dropping the constant $cK-\nu_{gt}$ of \eqref{eq:redcost}, that is
\begin{align}
    \max z = \sum_{c}\pi_{ct}\,a_{c}
    \label{eq:PPObjective}
\end{align}
subject to
\begin{align}
 \sum_{h\in H}\chi_{ih}=u_{i},\qquad \varrho_{ih}\le\chi_{ih}
   \quad & \forall i\le S,\;h\in H \label{eq:PPtype}\\
 \delta^{x}_{i}=\sum_{h}\bigl(p_{h}\chi_{ih}+(q_{h}-p_{h})\varrho_{ih}\bigr)
   \quad & \forall i\le S \label{eq:PPdimx}\\
 \delta^{y}_{i}=\sum_{h}\bigl(q_{h}\chi_{ih}+(p_{h}-q_{h})\varrho_{ih}\bigr)
   \quad & \forall i\le S \label{eq:PPdimy}\\
 \textsc{NoOverlap2D}\Bigl(\bigl\{[x_{i},x_{i}+\delta^{x}_{i})\times
   [y_{i},y_{i}+\delta^{y}_{i})\bigr\}_{i\,:\,u_{i}=1}\Bigr)
   \quad & \label{eq:PPnooverlap}\\
 x_{i}+\delta^{x}_{i}\le pK,\qquad y_{i}+\delta^{y}_{i}\le qK
   \quad & \forall i\le S \label{eq:PPcontain}\\
 \sum_{i\le S}\sum_{h} p_{h}q_{h}\chi_{ih}\;\le\;pK\cdot qK
   \quad & \label{eq:PParea}\\
 \ell_{ir}\;\le\;\bigl\lfloor rK/r_{gr}\bigr\rfloor
   \sum_{h:r\in R(h)}\chi_{ih}
   \quad & \forall i\le S,\;r\in R \label{eq:PPlayercompat}\\
 \sum_{r\in R} r_{gr}\,\ell_{ir}\;\le\;rK
   \quad & \forall i\le S \label{eq:PPheight}\\
 \sum_{c\,:\,\rho(c)=r} v_{ic}=n_{hr}\ell_{ir}
   \quad & \forall i\le S,\;r\in R \label{eq:PPassign}\\
 \sum_{c} w_{c}\,v_{ic}\;\le\;W^{\max}
   \quad & \forall i\le S \label{eq:PPcrush}\\
 a_{c}=\sum_{i\le S} v_{ic}
   \quad & \forall c \label{eq:PPagg}\\
 \sum_{c} w_{c}\,a_{c}\;\le\;pK^{w}
   \quad & \label{eq:PPpayload}\\
 \sum_{c}\gamma_{g\rho(c)}a_{c}\;\ge\;\alpha\,\Gamma_{K}
   \quad & \text{if }\theta=1 \label{eq:PPftlvol}\\
 \sum_{c} w_{c}\,a_{c}\;\ge\;\alpha\,pK^{w}
   \quad & \text{if }\theta=0 \label{eq:PPftlwgt}\\
 \sum_{i\le S}\sum_{h} p_{h}q_{h}\chi_{ih}\;\ge\;\alpha\,pK\cdot qK
   \quad & \text{if }\theta=1 \label{eq:PPareacut}\\
 \sum_{h}h\,\chi_{ih}\;\ge\;\sum_{h}h\,\chi_{i+1,h}
   \quad & \forall i<S. \label{eq:PPsymm}
\end{align}
Constraints \eqref{eq:PPtype}--\eqref{eq:PPagg} restrict the loads
used in the pricing problem to feasible ones, reproducing
\eqref{eq:col-geom}--\eqref{eq:col-agg}.
Equation \eqref{eq:PParea} is a redundant area cut, and
\eqref{eq:PPareacut} its counterpart under the volume regime: a slot
contributes at most (footprint area)$\times rK$ of RTI volume, so the
utilisation floor implies a floor on the occupied floor area. The
latter is what allows the solver to find a \emph{first} feasible
layout quickly, which for tight instances is harder than optimising
once one is known. Constraint \eqref{eq:PPsymm} breaks the symmetry
among interchangeable slots by ordering them on the stack-type index.

The subproblem is stated as a constraint program because its hard
core is the disjunctive non-overlap relation
\eqref{eq:PPnooverlap}, which a CP solver propagates natively through
a global constraint, whereas a MILP must encode it with the four
big-$M$ disjunctions
\eqref{sb:eq:nonov-x1}--\eqref{sb:eq:nonov-y2} and accept the weak
relaxation they entail. Section \ref{sec:results} reports the
resulting difference.

\paragraph{Making the oracle affordable.}
Pricing dominates the running time of the tree completely: on a
representative run it accounted for $1451$ of $1453$ seconds, against
two seconds for every master LP solved at every node, and a comparable
figure for the incumbent heuristic. Three measures bring it down
without changing what is solved. A node prices only the directions and
periods whose dispatch or delivery totals are fractional, those being
where the relaxation is interpolating between columns it already
holds. A node's target is derived from the master's residual, which
barely moves between a parent and its child, so a target already
priced on the same branch is skipped rather than solved again. And
each call receives a quarter of its budget first, the full budget only
if the short call returns nothing. Together these took a node from
about $28$ seconds to a few milliseconds, and raised the number of
nodes explored within a fixed budget by more than an order of
magnitude.

\subsection{Initial Columns}
\label{sec:initcols}

Column generation converges from any starting set, but a good one
shortens it and leaves the integer master more freedom to combine
loads. We initialise $\bar P^{g}$ by solving
\eqref{eq:PPObjective}--\eqref{eq:PPsymm} under a family of
demand-derived price vectors rather than duals: the unit vector,
which maximises the RTI count; every single-commodity vector
$\mathbf{e}_{c}$, which yields pure loads and keeps the master
flexible; the aggregate demand mix and each per-period demand mix;
and a set of random draws from the simplex for diversification.

A pool built from price vectors alone is \emph{extreme by
construction}. For any weights, \eqref{eq:PPObjective} attains its
optimum at an extreme point of $\operatorname{conv}(P^{g})$, so no
demand-derived vector --- aggregate, per-period, random or pure ---
can return a load from the interior of that hull. Yet interior mixes
are precisely what an integer plan needs when demand does not
decompose into full trucks, and on one instance the entire supplier
side of the pool consisted of two pure loads for this reason. We
therefore also seed by \emph{composition}: given a wanted mix
$\hat{a}$, the oracle minimises the volume-weighted asymmetric
deviation
$\sum_{c}u_{c}\bigl(\kappa\,\mathrm{under}_{c}+\mathrm{over}_{c}\bigr)$
with $\kappa>1$, subject to the same geometry, payload, crush and
full-truckload constraints. Because the target enters only the
objective, a mix too small to fill a truck degrades into the closest
fillable load rather than into infeasibility --- which is what the
last truck of a period needs. Seeding the per-period demand mixes this
way took that supplier side from two pure loads to five mixed ones.

Two devices make this robust. A constructive \emph{grid layout} ---
tile the floor with a single stack type in its better orientation and
fill each slot to its height, weight and crush limit --- is supplied
to the CP solver as a hint and retained as a fallback column. For the
near-tiling footprints of a real RTI catalogue it is already close to
the best attainable utilisation, and it guarantees that no direction
is ever left without columns. And any pricing call that returns
nothing is retried with a longer limit before being believed, since
an empty answer usually means the solver was starved of time rather
than that no load exists.

Before the integer master is solved, the pool is enriched with
\emph{reduced} loads obtained by unloading RTIs from existing
columns. Removing RTIs from a valid layout cannot break the geometry,
the height budget, a crush limit or the payload, so every reduced
load that still clears the full-truckload floor is itself feasible.
This matters because columns priced at the LP optimum are close to
full trucks, and integer combinations of full trucks are too coarse
to land inside the narrow delivery band left open by the supplier's
coverage stock. The LP hides this, because it can interpolate between
columns; the integer master cannot.

\subsection{Branching}
\label{sec:branching}

Branching on the master's own variables is the obvious rule and the
wrong one here. Fixing the $\lambda^{g}_{pt}$ with the largest
fractional part removes one load from the relaxation, but a
Dantzig--Wolfe master of this kind holds many columns that differ only
in detail, so the relaxation answers by shifting the same fractional
weight onto a near-identical column. Run to exhaustion on one
benchmark instance with the pool frozen, that rule branched $3000$
times and reached depth $1971$ without ever producing an integral
relaxation.

We therefore branch on aggregates first, taking the shallowest level
that is fractional:
\begin{enumerate}
\item the \emph{dispatch total} $\sum_{p}\lambda^{g}_{pt}$ of a
  direction and period. It is integer in every feasible plan and has
  no symmetric twin: every column of that period lies on one side of
  the row, so the dichotomy partitions plans rather than one
  variable's value.
\item the \emph{delivery totals} $Q_{nt}$ and $Q^{r}_{rt}$. These
  count physical units and are integer in any plan, although the
  relaxation leaves them continuous. They are original variables, so
  bounding one tightens a row that already carries a dual and the
  pricing subproblem keeps its structure.
\item the column variables $\lambda^{g}_{pt}$, as a last resort.
\end{enumerate}
The two aggregate levels absorb the fractionality almost entirely: on
a representative run, $61$ of $90$ branchings were on delivery totals,
$29$ on dispatch totals, and none on a column variable.

Both children are kept and the node with the cheapest relaxation is
expanded next, so a bad first move is recovered rather than carried to
the bottom. Every node re-prices, and that is what makes the tree
worth its cost. The pricing objective \eqref{eq:PPObjective} is
linear, so for \emph{any} dual vector its optimum is an extreme point
of $\operatorname{conv}(P^{g})$. A balanced load lying inside that
hull is dominated by a fractional combination of extreme loads and can
therefore never be priced, however the duals move --- the relaxation
buys the combination for one truck, while an integer plan must pay for
each of them. Branching down withdraws the dominating columns, the
duals move, and the oracle is asked for a load that does not yet
exist. On the hardest of our instances the search derived in this way
the only mixed supplier load a twelve-truck plan admits, which column
generation had not produced in $978$ columns; with it, a plan one
truck smaller became reachable.

Reaching integrality by branching alone remains slow, so the integer
master is solved over the current pool every $20$ nodes and its
solution kept as incumbent. It costs milliseconds at these pool sizes,
and without it the open list only grows, there being no incumbent
against which to prune.

\begin{algorithm}[!h]
\caption{Branch and price for 3D-SCPTOP}\label{alg:bp3d}
\begin{algorithmic}[1]
\State $\bar P \gets \textsc{InitialColumns}()$
       \Comment{Section \ref{sec:initcols}}
\State $sol_{int} \gets False$
\While{$sol_{int} == False$}
   \State $rc \gets -1$
   \While{$rc < 0$}
      \State $\pi,\nu,\lambda \gets \textsc{SolveRMP}(\bar P)$
      \State $rc, p \gets \textsc{Price}(\pi,\nu,\tau_{\mathrm{fast}})$
      \If{$rc \ge 0$}
         \State $rc, p \gets \textsc{Price}(\pi,\nu,\tau_{\mathrm{exact}})$
         \Comment{a fast pass may simply have missed a load}
      \EndIf
      \If{$rc < 0$}
         \State Append $p$ to $\bar P$
      \ElsIf{the pricing bound proves $rc \ge 0$}
         \State $\mathrm{LB} \gets$ RMP value \Comment{certified, Section \ref{sec:mp}}
      \EndIf
   \EndWhile
   \If{$\lambda$ is integral}
      \State $sol_{int} \gets True$; \textbf{break}
   \EndIf
   \State $(p,t) \gets \arg\max$ fractional part of $\lambda^{g}_{pt}$
   \State fix $\lambda^{g}_{pt} \ge \lceil \lambda^{g}_{pt}\rceil$,
          or $\le \lfloor \lambda^{g}_{pt}\rfloor$ if that node is infeasible
\EndWhile
\State $\bar P \gets \bar P \cup \textsc{ReducedLoads}(\bar P)$
\State $\mathrm{UB} \gets \textsc{SolveRMP}(\bar P)$ as an integer program
\State verify every dispatched column against
       \eqref{eq:col-geom}--\eqref{eq:col-agg}
\State \Return $(\mathrm{UB}, \mathrm{LB})$
\end{algorithmic}
\end{algorithm}


\subsection{Two corrections to the compact model}
\label{sec:corrections}

Implementing 3D-SCPTOP surfaced two defects worth recording, both
invisible at $lt=1$.

\paragraph{Inventory recursion under a lead time.}
Constraints \eqref{sb:eq:inv-plant-main} and
\eqref{sb:eq:inve-sup-main} carry inventory forward from period
$t-lt$, whereas inventory accumulates one period at a time and only
the \emph{arrival} is delayed. They must read
\begin{align}
I_{pnt}  &=I_{p,n,t-1}-d_{pnt}+Q_{n,t-lt}
    &&\forall n\in N,\; t\ge lt+1,
    \label{eq:fix-plant}\\
Ie_{str} &=Ie_{s,t-1,r}+Q_{r,t-lt}-\textstyle\sum_{n\in M(r)} d_{snt}
    &&\forall r\in R,\; t\ge lt+1.
    \label{eq:fix-sup}
\end{align}
With $lt=1$ the two versions coincide, which is why the defect
survived; with $lt\ge2$ the original silently discards $lt-1$ periods
of stock.

\paragraph{Coverage stock on folded RTIs.}
Requiring $Ie_{gtr}\ge ss_{gtr}$ at \emph{both} nodes, as
\eqref{sb:eq:safety} does, forces the plant to hoard empties it has
no use for and can render an otherwise feasible instance infeasible
when the RTI fleet is tight. The operationally meaningful requirement
is that the \emph{supplier} hold enough folded RTIs to cover its
production, so we impose the bound at $g=s$ only.

\endgroup
